\chapter{Em quên vở rồi ạ}
\begin{tinhhuongbox}{Tình huống | Situation}
\songngu{Một học sinh quên vở đầu tiết, và nếu thầy cô đẩy quá mạnh, cả lớp sẽ mở bài bằng căng thẳng.}{A student forgets their notebook at the start, and if you push too hard, the whole class will open with tension.}
\songngu{Lúc này cần giữ mặt lớp mà vẫn giữ việc học.}{At this moment, you need to save the class's face while still saving the lesson.}
\end{tinhhuongbox}
\begin{noitrenlop}
\songngu{Không sao, hôm nay em nhớ bài bằng đầu nhiều hơn bằng giấy nhé.}{That is okay. Today you will remember the lesson more with your head than with paper.}
\end{noitrenlop}
\begin{vihieuqua}
\songngu{Câu này hạ xung đột nhưng không hạ trách nhiệm.}{This line lowers the conflict without lowering responsibility.}
\end{vihieuqua}
\begin{englishpocket}
\songngu{Cụm nên nhớ: \textit{That is okay.}}{Phrase to keep: \textit{That is okay.}}
\end{englishpocket}
\chapter{Em đi học muộn ạ}
\begin{tinhhuongbox}{Tình huống | Situation}
\songngu{Đi muộn mà xử lý dài dòng thì cái muộn còn kéo dài thêm.}{If you handle lateness with a long scene, the lateness itself gets longer.}
\songngu{Cần nhắc ngắn rồi đưa học sinh về bài ngay.}{You need a short reminder and then a quick return to the lesson.}
\end{tinhhuongbox}
\begin{noitrenlop}
\songngu{Em đến hơi trễ rồi, giờ đừng trễ luôn cả sự tập trung nữa nhé.}{You are a little late already, so do not let your focus arrive late too.}
\end{noitrenlop}
\begin{vihieuqua}
\songngu{Câu này chạm đúng việc cần làm ngay sau khi ngồi xuống.}{This line touches the exact thing the student must do right after sitting down.}
\end{vihieuqua}
\begin{englishpocket}
\songngu{Cụm nên nhớ: \textit{Do not let your focus arrive late.}}{Phrase to keep: \textit{Do not let your focus arrive late.}}
\end{englishpocket}
\chapter{Bạn kia cũng chưa làm mà}
\begin{tinhhuongbox}{Tình huống | Situation}
\songngu{Khi học sinh kéo người khác vào làm lá chắn, nhịp lớp sẽ bắt đầu trôi sang chuyện khác.}{When a student pulls someone else in as a shield, the class rhythm starts drifting into something else.}
\songngu{Phải kéo trọng tâm về đúng người đang nói chuyện với mình.}{You have to pull the focus back to the person you are actually talking to.}
\end{tinhhuongbox}
\begin{noitrenlop}
\songngu{Bạn kia lát nữa sẽ có phần của bạn ấy, còn bây giờ tôi đang nói chuyện với em.}{That student will have their turn later, but right now I am talking to you.}
\end{noitrenlop}
\begin{vihieuqua}
\songngu{Câu này bình tĩnh mà rất chắc tay.}{This line is calm, but it is very firm.}
\end{vihieuqua}
\begin{englishpocket}
\songngu{Cụm nên nhớ: \textit{Right now I am talking to you.}}{Phrase to keep: \textit{Right now I am talking to you.}}
\end{englishpocket}
\chapter{Em chỉ đùa thôi ạ}
\begin{tinhhuongbox}{Tình huống | Situation}
\songngu{Có những câu đùa không ác nhưng kéo nhịp tiết đi lạc.}{Some jokes are not mean, but they still pull the lesson off track.}
\songngu{Lúc đó nên chỉnh hướng tiếng cười, không nên bóp chết nó hoàn toàn.}{At that point, you should redirect the laughter instead of killing it completely.}
\end{tinhhuongbox}
\begin{noitrenlop}
\songngu{Câu đó giữ lại cho giờ ra chơi nhé, còn giờ cho bài học vui ké với.}{Save that line for recess, and for now let the lesson have some fun too.}
\end{noitrenlop}
\begin{vihieuqua}
\songngu{Học sinh thấy mình không bị dập, chỉ bị kéo về đúng sân chơi.}{Students feel that they are not being shut down, only brought back to the right stage.}
\end{vihieuqua}
\begin{englishpocket}
\songngu{Cụm nên nhớ: \textit{Save that line for recess.}}{Phrase to keep: \textit{Save that line for recess.}}
\end{englishpocket}
\chapter{Lớp đang ồn quá}
\begin{tinhhuongbox}{Tình huống | Situation}
\songngu{Lớp ồn vì nhiều em cùng muốn nói, chứ không hẳn vì chống đối.}{The class is noisy because many students want to speak at once, not necessarily because they are resisting.}
\songngu{Thầy cô cần hạ âm lượng mà không hạ ý kiến.}{You need to lower the volume without lowering the ideas.}
\end{tinhhuongbox}
\begin{noitrenlop}
\songngu{Mình hạ âm lượng chứ không hạ chất lượng ý kiến nhé.}{Let us lower the volume, not the quality of the ideas.}
\end{noitrenlop}
\begin{vihieuqua}
\songngu{Câu này khiến lớp thấy mình đang được chỉnh cách nói, không phải bị cấm nói.}{This line makes the class feel that their way of speaking is being adjusted, not their right to speak being removed.}
\end{vihieuqua}
\begin{englishpocket}
\songngu{Cụm nên nhớ: \textit{Lower the volume, not the quality.}}{Phrase to keep: \textit{Lower the volume, not the quality.}}
\end{englishpocket}